\documentclass[10pt,a4paper]{article}
\usepackage[margin=2.2cm]{geometry}
\usepackage[T1]{fontenc}
\usepackage{lmodern}
\usepackage{booktabs}
\usepackage{tabularx}
\usepackage{enumitem}
\usepackage{xcolor}
\usepackage[hidelinks]{hyperref}
\setlist{itemsep=2pt,topsep=3pt}
\setlength{\parindent}{0pt}
\setlength{\parskip}{5pt}
\newcommand{\val}[1]{\texttt{#1}}
\newcolumntype{L}{>{\raggedright\arraybackslash}X}

\title{GHAgentCI validation: coding guide}
\author{Precision and recall labelling, MSR 2027}
\date{Version 2, 29 October 2026}

\begin{document}
\maketitle
\vspace{-1.5em}

\section*{In one paragraph}

\textsf{awe} scans GitHub Actions workflow files and reports the places where
it believes a \emph{coding agent} (Claude, Codex, Copilot, \dots) is started by
CI. You will check a sample of those reports, and a sample of workflows it did
\emph{not} flag, by opening each workflow on GitHub and answering two
questions. Your answers measure how often \textsf{awe} is right (precision) and
what it misses (recall). Two coders answer independently; agreement between you
is part of the result.

\section*{1. What you received}

\begin{tabularx}{\linewidth}{@{}l L@{}}
\toprule
\textbf{File} & \textbf{What it is} \\
\midrule
\val{GHAgentCI-labeling-coder-X.xlsx} & Your workbook. Tabs: \emph{Start here}, \emph{Codebook} (quick reference), \emph{Precision} (576 units), \emph{Recall} (200 units). Only yellow cells are yours to fill. \\
\val{GHAgentCI-coding-guide.pdf} & This guide. \\
\val{GHAgentCI-calibration.xlsx} & (shared folder) 30 practice units, coded \emph{together} in the calibration session. Not counted. \\
\bottomrule
\end{tabularx}

\section*{2. Steps and timing}

\begin{enumerate}
\item \textbf{Calibration session (together, about 1 hour).} Code the 30 units of
\val{GHAgentCI-calibration.xlsx} jointly, discuss every disagreement, and settle
the open rule in Section~6. This guide is updated afterwards if a rule changes.
\item \textbf{Precision tab, alone (about 10 hours).} 576 units.
\item \textbf{Recall tab, alone (about 3 hours).} 200 units.
\item \textbf{Return your workbook by 3 November 2026}, unchanged in structure, to
Daniel Ogenrwot. Disagreements are then resolved together on 4 November.
\end{enumerate}

At about one minute per unit, that is roughly 13 hours of coding: plan for about
two and a half hours a day until 3 November. Save often.

\section*{3. Rules}

\begin{itemize}
\item \textbf{Work independently.} Do not discuss units with the other coder,
and do not look at their workbook, until both are returned. Ask questions about
the \emph{guide} to Daniel instead.
\item \textbf{Code it yourself.} Do not use AI assistants or scripts to answer;
the point is a human judgement.
\item \textbf{Judge the workflow, not the detector.} The workbook deliberately
does not show what \textsf{awe} decided. Do not look it up.
\item \textbf{Do not add, delete, rename or reorder columns.} Sorting or
filtering rows is fine; every answer is matched by \val{unit\_id}.
\item \textbf{When in doubt, write a short note} in the \val{notes} column. Notes
are read during adjudication and cost you nothing.
\end{itemize}

\section*{4. How to look at a unit}

Click the link in the \val{open workflow} column. It opens the workflow file on
GitHub \emph{at the commit being judged}, not today's version.

\begin{itemize}
\item \val{scope = step}: judge the step identified by \val{job\_id} and
\val{step\_index}. Steps are counted from \val{0} within their job: index
\val{0} is the first entry under that job's \val{steps:}. \val{step\_name}, when
present, is the step's \val{name:}. An index of \val{-1} means the job itself
calls another workflow (\val{jobs.<job>.uses:}); judge that call.
\item \val{scope = file}: judge the workflow as a whole. This is used when the
evidence concerns the entire file: a reserved file name, a credential set for
the whole workflow, or a GitHub Agentic Workflows (gh-aw) file.
\item \textbf{Recall tab:} units have no step. Judge the whole file: does it
invoke a coding agent \emph{anywhere}?
\item If a link no longer resolves (repository deleted or made private), answer
\val{maybe} for Q1 and write \val{link dead} in \val{notes}.
\end{itemize}

\section*{5. Q1 --- Does this step (or file) invoke a coding agent when the workflow runs?}

A \emph{coding agent} is an LLM-based tool that acts on the repository or its
activity: it reads code, diffs, issues or comments, and produces reviews,
comments, labels, commits or pull requests.

\begin{tabularx}{\linewidth}{@{}l L@{}}
\toprule
\textbf{Answer} & \textbf{Use when} \\
\midrule
\val{yes} & The step (or file) runs such an agent when the workflow executes. \\
\val{maybe} & An LLM is called but it is unclear that it acts as an agent, or the agent step cannot actually run (disabled, dry-run only). \\
\val{no} & No agent runs. The match is incidental. \\
\bottomrule
\end{tabularx}

\subsection*{Decision rules for common cases}

\begin{tabularx}{\linewidth}{@{}L l@{}}
\toprule
\textbf{Case} & \textbf{Q1} \\
\midrule
Step uses a vendor's agent action, e.g.\ \val{anthropics/claude-code-action}, \val{openai/codex-action}, \val{google-github-actions/run-gemini-cli} & \val{yes} \\
Shell step runs an agent CLI, e.g.\ \val{claude -p "\dots"}, \val{codex exec}, \val{gemini -p}, \val{aider -{}-message} & \val{yes} \\
A third-party action that wraps an agent (e.g.\ an "AI PR reviewer" calling Claude or GPT on the diff and posting a review) & \val{yes} \\
File \val{.github/workflows/copilot-setup-steps.yml} (prepares the Copilot coding agent's environment) & \val{yes} \\
gh-aw file: a compiled \val{*.lock.yml} with a "generated by gh-aw" header, or its Markdown source with gh-aw frontmatter & \val{yes} \\
Agent step guarded by \val{if: false}, or a mode that only prints what it would do & \val{maybe} \\
Single LLM call that only summarises or formats text (release notes, changelog), not acting on code & \val{maybe} \\
Step only \emph{installs} an agent (\val{npm i -g @openai/codex}) and nothing in that step runs it & \val{no} \\
An API key such as \val{ANTHROPIC\_API\_KEY} feeds a test suite or an application, not an agent & \val{no} \\
The agent's name appears only in an \val{echo}, a comment, documentation, or a file path & \val{no} \\
An unrelated tool that shares an agent's name & \val{no} \\
File is not a working workflow (e.g.\ ordinary Markdown in \val{.github/workflows/}) & \val{no} \\
Invalid YAML: judge what the author evidently intended the step to do & as above \\
\bottomrule
\end{tabularx}

\section*{6. Settled at calibration}

\textbf{Open until the calibration session:} a step that only calls a generic
LLM inference action, such as \val{actions/ai-inference}, to summarise or
classify an issue. Is that \val{yes} or \val{maybe}?

\medskip
Decision: \rule{0.6\linewidth}{0.4pt}

\section*{7. Q2 --- Which agent is it?}

Answer Q2 only when Q1 is \val{yes} or \val{maybe}; leave it empty for \val{no}.
Pick one value from the dropdown. Spelling matters: answers are compared
literally.

\begin{tabularx}{\linewidth}{@{}l L@{}}
\toprule
\textbf{Value} & \textbf{Use when} \\
\midrule
\val{claude} \val{codex} \val{copilot} \val{gemini} & The vendor's agent, directly or through a wrapper action that clearly uses it. \\
\val{cursor} \val{aider} \val{devin} \val{pr\_agent} & \val{pr\_agent} is Qodo's PR-Agent (\val{Codium-ai/pr-agent}). \\
\val{coderabbit} \val{goose} \val{opencode} \val{openhands} & \\
\val{generic\_ai} & A third-party AI action not tied to one of the vendors above. \\
\val{unknown} & A gh-aw file that does not say which engine it uses. \\
\val{other} & A coding agent not in this list. \\
\val{unclear} & Cannot be determined from the workflow. \\
\bottomrule
\end{tabularx}

For gh-aw files, the engine is the \val{engine:} key of the Markdown source, or
the \val{agent\_id} in the lock file's \val{gh-aw-metadata} header; if neither
is present, answer \val{unknown}.

\section*{8. Worked examples}

\begin{tabularx}{\linewidth}{@{}l l L l l@{}}
\toprule
\textbf{scope} & \textbf{job / step} & \textbf{What you see} & \textbf{Q1} & \textbf{Q2} \\
\midrule
step & \val{review} / 1 & \val{uses: anthropics/claude-code-action@v1} with a review prompt & \val{yes} & \val{claude} \\
step & \val{test} / 3 & \val{run: pytest}, with \val{OPENAI\_API\_KEY} in \val{env:} & \val{no} & \\
file & -- & \val{copilot-setup-steps.yml} installing project dependencies & \val{yes} & \val{copilot} \\
step & \val{triage} / 0 & \val{run: echo "ask @claude for help"} & \val{no} & \\
file & -- & \val{issue-triage.lock.yml} whose header names engine \val{codex} & \val{yes} & \val{codex} \\
\bottomrule
\end{tabularx}

\section*{9. Questions}

Ask Daniel Ogenrwot (\href{mailto:ogenrwot@unlv.nevada.edu}{ogenrwot@unlv.nevada.edu}).
Questions about a specific unit go in its \val{notes} instead, so that they do not
influence the other coder.

\end{document}
